% Template para uso no TCC do Departamento de Física de Ji-Paraná
% da Universidade Federal de Rondônia
% Elaborado em 10 de junho de 2017 por Marco Polo Moreno de Souza com base no template book.
% Aperfeiçoado em 07 de março de 2019, com a introdução de macros
% Ajustado em 23 de junho de 2023 (folha de rosto + ficha catalográfica; remoção da ata de aprovação)

\documentclass[a4paper,reqno,12pt,oneside]{book}

\usepackage{amsmath}%
\usepackage{amsfonts}%
\usepackage{amssymb}%
\usepackage{graphicx}

\usepackage[notlof,notlot,nottoc]{tocbibind}
\usepackage{glossaries}
\makenoidxglossaries
\glstoctrue
\input{macros/glossary.tex}

%Veja o arquivo abaixo antes de adicionar outras bibliotecas
\input{macros/bibliotecas.tex}


%%%%%%%%%%%%%%%%%%%%%%%%%%%%%%%%%%%%%%%%%%%%%%%%%%%%%%%%%%%%%%%
%Preencha os campos abaixo conforme seu TCC
\newcommand{\Titulo}{Título do trabalho, seguido do subtítulo, se houver}
\newcommand{\Nome}{Seu nome}
\newcommand{\curso}{Licenciatura Plena} %Digite Licenciatura Plena ou Bacharelado
\newcommand{\graduacao}{Licenciado} %Digite Licenciado ou Bacharel conforme o caso
\newcommand{\orientador}{Nome do orientador} %Nome com título. Ex: Dr. Marco Polo Moreno de Souza
\newcommand{\primeiromembro}{Nome do professor da banca} %Nome com título. Ex: Dr. Fulano de Tal
\newcommand{\segundomembro}{Nome do professor da banca} %Nome com título. Ex: MSc. Beltrano
\newcommand{\instituicao}{xxxxx} %Instituição do segundo membro. Se for interno, digite DEFIJI/CJP/UNIR
\newcommand{\Cidade}{Ji-Paraná, RO}
\newcommand{\Data}{Mês e ano da defesa} %Ex: Junho de 2020


%%%%%%%%%%%%%%%%%%%%%%%%%%%%%%%%%%%%%%%%%%%%%%%%%%%%%%%%%%%%%%%
%Dados da defesa, a ser preenchido pelo orientador
\newcommand{\dia}{xxx} %Ex: quinze
\newcommand{\mes}{xxx} %Ex: abril
\newcommand{\ano}{xxx} %Ex: 2019
\newcommand{\inicio}{xxx} %Início da defesa. Ex: 10:15 h
\newcommand{\fim}{xxx} %Fim da sessão. Ex: 11:45 h
\newcommand{\arguicao}{xxx} %Duração da arguição da banca em minutos. Ex: 45
\newcommand{\arguicaoextenso}{xxx} %Idem acima, por extenso. Ex: quarenta e cinco
\newcommand{\local}{xxx} %Ex: auditório do Campus da UNIR de Ji-Paraná
\newcommand{\nota}{xxx} %Ex: 8,5
\newcommand{\notaextenso}{xxx} %Ex: oito inteiros e cinco décimos

\input{macros/macros.tex}
\input{macros/estilo.tex}

\begin{document}
\renewcommand{\bibname}{Referências}

\input{macros/capa.tex}
\input{macros/contra-capa.tex}
\includepdf[pages=1,width=1.2\textwidth, offset=0 1cm]{anexos/ficha-catalografica.pdf}
\includepdf[pages={1,2},width=1.0\textwidth, offset=0 -2cm]{anexos/folha-rosto-sei.pdf}
%\input{macros/ata.tex}

%%%%%%%%%%%%%%%%%%%%%%%%%%%%%%%%%%%%%%%%%%%%%%%%%%%%%%%%%%%%%%%
\chapter*{Dedicatória}
\thispagestyle{empty}
%A organização das páginas de dedicatória, agradecimentos e epígrafe ficam a critério do aluno, sendo opcionais.

Digite a dedicatória aqui.


%%%%%%%%%%%%%%%%%%%%%%%%%%%%%%%%%%%%%%%%%%%%%%%%%%%%%%%%%%%%%%%
\chapter*{Agradecimentos}
\thispagestyle{empty}

Digite os agradecimentos aqui.


%%%%%%%%%%%%%%%%%%%%%%%%%%%%%%%%%%%%%%%%%%%%%%%%%%%%%%%%%%%%%%%
\chapter*{Epígrafe}
\thispagestyle{empty}

Digite a epígrafe aqui.


%%%%%%%%%%%%%%%%%%%%%%%%%%%%%%%%%%%%%%%%%%%%%%%%%%%%%%%%%%%%%%%
\chapter*{Resumo}
\thispagestyle{empty}
\noindent
O resumo em língua portuguesa deverá conter no mínimo 150 e no máximo 500 palavras.
Bla Bla Bla Bla Bla Bla Bla Bla Bla Bla Bla Bla Bla Bla Bla Bla Bla Bla 
Bla Bla Bla Bla Bla Bla Bla Bla Bla Bla Bla Bla Bla Bla Bla Bla Bla Bla 
Bla Bla Bla Bla Bla Bla Bla Bla Bla Bla Bla Bla Bla Bla Bla Bla Bla Bla 
Bla Bla Bla Bla Bla Bla Bla Bla Bla Bla Bla Bla Bla Bla Bla Bla Bla Bla 
Bla Bla Bla Bla Bla Bla Bla Bla Bla Bla Bla Bla Bla Bla Bla Bla Bla Bla 
Bla Bla Bla Bla Bla Bla Bla Bla Bla Bla Bla Bla Bla Bla Bla Bla Bla Bla 
Bla Bla Bla Bla Bla Bla Bla Bla Bla Bla Bla Bla Bla Bla Bla Bla Bla Bla 
Bla Bla Bla Bla Bla Bla Bla Bla Bla Bla Bla Bla Bla Bla Bla Bla Bla Bla 
Bla Bla Bla Bla Bla Bla Bla Bla Bla Bla Bla Bla Bla Bla Bla Bla Bla Bla 
Bla Bla Bla Bla Bla Bla Bla Bla Bla Bla Bla Bla Bla Bla Bla Bla Bla Bla 
Bla Bla Bla Bla Bla Bla Bla Bla Bla Bla Bla Bla Bla Bla Bla Bla Bla Bla 
Bla Bla Bla Bla Bla Bla Bla Bla Bla Bla Bla Bla Bla Bla Bla Bla Bla Bla 
Bla Bla Bla Bla Bla Bla Bla Bla Bla Bla Bla Bla Bla Bla Bla Bla Bla Bla 
Bla Bla Bla Bla Bla Bla Bla Bla Bla Bla Bla Bla Bla Bla Bla Bla Bla Bla 
Bla Bla Bla Bla Bla Bla Bla Bla Bla Bla Bla Bla Bla Bla Bla Bla Bla Bla 
Bla Bla Bla Bla Bla Bla Bla Bla Bla Bla Bla Bla Bla Bla Bla Bla Bla Bla 
Bla Bla Bla Bla Bla Bla Bla Bla Bla Bla Bla Bla Bla Bla Bla Bla Bla Bla 
Bla Bla Bla Bla Bla Bla Bla Bla Bla Bla Bla Bla Bla Bla Bla Bla Bla Bla 

\vspace{1.0cm}
\noindent
\textbf{Palavras-chave}: palavra-chave 1. palavra-chave 2. palavra-chave 3.
% Minimo de 3 e máximo de 6 palavras-chave.


%%%%%%%%%%%%%%%%%%%%%%%%%%%%%%%%%%%%%%%%%%%%%%%%%%%%%%%%%%%%%%%
\chapter*{Abstract}
\thispagestyle{empty}
\noindent
Type de abstract.
Bla Bla Bla Bla Bla Bla Bla Bla Bla Bla Bla Bla Bla Bla Bla Bla Bla Bla 
Bla Bla Bla Bla Bla Bla Bla Bla Bla Bla Bla Bla Bla Bla Bla Bla Bla Bla 
Bla Bla Bla Bla Bla Bla Bla Bla Bla Bla Bla Bla Bla Bla Bla Bla Bla Bla 
Bla Bla Bla Bla Bla Bla Bla Bla Bla Bla Bla Bla Bla Bla Bla Bla Bla Bla 
Bla Bla Bla Bla Bla Bla Bla Bla Bla Bla Bla Bla Bla Bla Bla Bla Bla Bla 
Bla Bla Bla Bla Bla Bla Bla Bla Bla Bla Bla Bla Bla Bla Bla Bla Bla Bla 
Bla Bla Bla Bla Bla Bla Bla Bla Bla Bla Bla Bla Bla Bla Bla Bla Bla Bla 
Bla Bla Bla Bla Bla Bla Bla Bla Bla Bla Bla Bla Bla Bla Bla Bla Bla Bla 
Bla Bla Bla Bla Bla Bla Bla Bla Bla Bla Bla Bla Bla Bla Bla Bla Bla Bla 
Bla Bla Bla Bla Bla Bla Bla Bla Bla Bla Bla Bla Bla Bla Bla Bla Bla Bla 
Bla Bla Bla Bla Bla Bla Bla Bla Bla Bla Bla Bla Bla Bla Bla Bla Bla Bla 
Bla Bla Bla Bla Bla Bla Bla Bla Bla Bla Bla Bla Bla Bla Bla Bla Bla Bla 
Bla Bla Bla Bla Bla Bla Bla Bla Bla Bla Bla Bla Bla Bla Bla Bla Bla Bla 
Bla Bla Bla Bla Bla Bla Bla Bla Bla Bla Bla Bla Bla Bla Bla Bla Bla Bla 
Bla Bla Bla Bla Bla Bla Bla Bla Bla Bla Bla Bla Bla Bla Bla Bla Bla Bla 
Bla Bla Bla Bla Bla Bla Bla Bla Bla Bla Bla Bla Bla Bla Bla Bla Bla Bla 
Bla Bla Bla Bla Bla Bla Bla Bla Bla Bla Bla Bla Bla Bla Bla Bla Bla Bla 
Bla Bla Bla Bla Bla Bla Bla Bla Bla Bla Bla Bla Bla Bla Bla Bla Bla Bla 

\vspace{1.0cm}
\noindent
\textbf{Keywords}: keyword 1. keyword 2. keyword 3.
% Minimo de 3 e máximo de 6 palavras-chave em inglês.

\input{macros/listas-sumario.tex}


%%%%%%%%%%%%%%%%%%%%%%%%%%%%%%%%%%%%%%%%%%%%%%%%%%%%%%%%%%%%%%%
\chapter{Introdução}
\label{introducao}

Digite a introdução aqui.


%%%%%%%%%%%%%%%%%%%%%%%%%%%%%%%%%%%%%%%%%%%%%%%%%%%%%%%%%%%%%%%
\chapter{Título do Segundo Capítulo}
\label{capitulo2}
%Para cada capítulo adicional, digite o comando \chapter{Título do Capítulo}

Digite o conteúdo do \footnote{Exemplo de nota de rodapé.}segundo capítulo aqui.

\section{Título de uma seção}
\label{secao2.1}

Digite o conteúdo do uma seção aqui.

\subsection{Subseção}

Digite o conteúdo do uma subseção aqui.

\section{Figuras}

Exemplo de figura:

\begin{figure}[H]
	\centering
		\includegraphics[width=0.60\textwidth]{figuras/fig1.eps}
	\label{fig1}
	\caption{Espectro de um laser de femtossegundos.}
\end{figure}

%Citação de figura: \ref{fig1}.

\section{Equações}

Exemplo de equação centralizada:

\begin{equation}
\label{pitagoras}
	a^2 = b^2 + c^2.
\end{equation}

Exemplo de equação no texto: $e^{ix} = \cos{x} + i\sin{x}$. Citação de equação: \ref{pitagoras}.

\section{Tabelas}

Exemplo de tabela:

\begin{table}[H]
\begin{center}
	\begin{tabular}{|c|c|c|}
		\hline
		$\sin{x}$ & $\cos{x}$ & $\tan{x}$ \\
		\hline
		$\sec{x}$ & $\csc{x}$ & $\cot{x}$ \\
		\hline
		$\arcsin{x}$ & $\arccos{x}$ & $\arctan{x}$ \\
		\hline
		$\sinh{x}$ & $\cosh{x}$ & $\tanh{x}$ \\
		\hline
	\end{tabular}
	\vspace{0.3cm}
	\caption{Funções trigonométricas e hiperbólicas.}
	\label{funcoes}
\end{center}
\end{table}

%Citação de tabela: \ref{funcoes}

\section{Códigos}

Exemplo de código (linguagem C):

\vspace{0.5cm}
\begin{lstlisting}
#include<stdio.h>

int k;

main()
{
	for (k=1; k<=5; k++)
		printf("Fisica - UNIR - Ji-Parana\n");
}
\end{lstlisting}

\section{Citação}

Exemplo de citação:

\vspace{0.5cm}
\begin{addmargin}[3cm]{0.0cm}\small{
		Bla Bla Bla Bla Bla Bla Bla Bla Bla Bla Bla Bla Bla Bla Bla Bla Bla Bla Bla Bla
		Bla Bla Bla Bla Bla Bla Bla Bla Bla Bla Bla Bla Bla Bla Bla Bla Bla Bla Bla Bla
		Bla Bla Bla Bla Bla Bla Bla Bla Bla Bla Bla Bla Bla Bla Bla Bla Bla Bla Bla Bla
		Bla Bla Bla Bla Bla Bla Bla Bla Bla Bla Bla Bla Bla Bla Bla Bla Bla Bla Bla Bla
		Bla Bla Bla Bla Bla Bla Bla Bla Bla Bla Bla Bla Bla Bla Bla Bla Bla Bla Bla Bla
		Bla Bla Bla Bla Bla Bla Bla Bla Bla Bla Bla Bla Bla Bla Bla Bla Bla Bla Bla Bla
		Bla Bla Bla Bla Bla Bla Bla Bla Bla Bla Bla Bla Bla Bla Bla Bla Bla Bla Bla Bla
}\end{addmargin}
\vspace{0.5cm}

Citando um trabalho: \cite{Nussenzveig} ou (Nussenzveig, 2015).


%%%%%%%%%%%%%%%%%%%%%%%%%%%%%%%%%%%%%%%%%%%%%%%%%%%%%%%%%%%%%%%
\chapter[Outro capítulo]{Outro capítulo}

Digite aqui o conteúdo de outro capítulo.

\begin{table}[H]
\begin{center}
	\begin{tabular}{|c|c|c|}
		\hline
		$\sin{x}$ & $\cos{x}$ & $\tan{x}$ \\
		\hline
		$\sec{x}$ & $\csc{x}$ & $\cot{x}$ \\
		\hline
		$\arcsin{x}$ & $\arccos{x}$ & $\arctan{x}$ \\
		\hline
		$\sinh{x}$ & $\cosh{x}$ & $\tanh{x}$ \\
		\hline
	\end{tabular}
	\vspace{0.3cm}
	\caption{Funções trigonométricas e hiperbólicas.}
	\label{funcoes2}
\end{center}
\end{table}

\begin{figure}[H]
	\centering
		\includegraphics[width=0.60\textwidth]{figuras/fig1.eps}
	\label{fig2}
	\caption{Espectro de um laser de femtossegundos.}
\end{figure}


%%%%%%%%%%%%%%%%%%%%%%%%%%%%%%%%%%%%%%%%%%%%%%%%%%%%%%%%%%%%%%%
\chapter{Conclusão}

Digite a conclusão do TCC aqui.


%%%%%%%%%%%%%%%%%%%%%%%%%%%%%%%%%%%%%%%%%%%%%%%%%%%%%%%%%%%%%%%
%Digite as referências abaixo
\begin{thebibliography}{99}

\bibitem {Nussenzveig} Herch Moysés Nussenzveig. \textbf{Curso de Física Básica: Eletromagnetismo}, volume 3 (Editora Edgard Blucher, São Paulo, 2015).

\bibitem {Abbott} B. P. Abbott \textit{et al}. \textbf{Observation of Gravitational Waves from a Binary Black Hole Merger}. \textit{Phys. Rev. Lett.} \textbf{116}, 061102 (2016).

\bibitem {Ross} I. M. Ross. \textit{et al}. \textbf{The Invention of the Transistor}. \textit{Proceedings of the IEEE} \textbf{86}, 7 (1998).
%\vspace{0.4cm}

\bibitem {Strickland} Donna Strickland -- Nobel Lecture. \url{https://www.nobelprize.org/prizes/physics/2018/strickland/lecture/} (acesso em 11 de março de 2019).

\end{thebibliography}

\backmatter \appendix


%%%%%%%%%%%%%%%%%%%%%%%%%%%%%%%%%%%%%%%%%%%%%%%%%%%%%%%%%%%%%%%
\chapter{Título do Primeiro Apêndice}

Digite o primeiro apêndice aqui.


%%%%%%%%%%%%%%%%%%%%%%%%%%%%%%%%%%%%%%%%%%%%%%%%%%%%%%%%%%%%%%%
\chapter{Título do Segundo Apêndice}
%Para cada apêndice adicional, insira o comando \chapter{Título do Apêndice}

Digite o segundo apêndice aqui.


%%%%%%%%%%%%%%%%%%%%%%%%%%%%%%%%%%%%%%%%%%%%%%%%%%%%%%%%%%%%%%%
\includepdf[scale=0.95, pages=1, pagecommand=\chapter{Anexo A -- Título do Anexo A}\label{anexo-a}, offset= 0.0cm -1.0cm]{anexos/anexo1.pdf}
%Para cada anexo adicional, insira o comando conforme acima


%%%%%%%%%%%%%%%%%%%%%%%%%%%%%%%%%%%%%%%%%%%%%%%%%%%%%%%%%%%%%%%
\includepdf[pages=1,width=1.3\textwidth, offset=0 -2cm]{anexos/termo-de-autorizacao.pdf}
%\thispagestyle{empty}

\printnoidxglossaries

\end{document}